% Generated draft: requires mathematical and source review.
\documentclass[11pt]{article}
\usepackage[margin=25mm]{geometry}
\usepackage{fontspec}
\setmainfont{Latin Modern Roman}
\usepackage{amsmath,amssymb,mathtools,unicode-math}
\setmathfont{Latin Modern Math}
\usepackage{enumitem,needspace,xurl,hyperref}
\setlength{\parindent}{0pt}
\setlength{\parskip}{5pt}
\setlength{\emergencystretch}{3em}
\newcommand{\N}{\mathbb{N}}
\newcommand{\Z}{\mathbb{Z}}
\newcommand{\Q}{\mathbb{Q}}
\newcommand{\R}{\mathbb{R}}
\newcommand{\C}{\mathbb{C}}
\newcommand{\F}{\mathbb{F}}
\newcommand{\A}{\mathbb{A}}
\newcommand{\PP}{\mathbb{P}}
\newcommand{\E}{\mathbb{E}}
\newcommand{\eps}{\varepsilon}
\newcommand{\dd}{\,\mathrm{d}}
\begin{document}
\section*{UnsolvedMath ID: 30000153}
\textbf{Source problem code:} catalogue entry
\subsection{Definitions and mathematical statement}
Consider an atom with nucleus at the origin and an electronic ground state in the symmetry subspace intended by the source. The Hamiltonian, electron number, nuclear charge, and symmetry sector are not specified in the extracted record. For a normalized \(N\)-electron wave function \(\Psi\), the usual physical one-particle density is
\[
\rho_{\mathrm{phys}}(x)=\sum_{j=1}^N\int |\Psi(x_1,\ldots,x_{j-1},x,x_{j+1},\ldots,x_N)|^2\prod_{\ell\ne j}\dd x_\ell,
\]
with a spin sum if spin variables are present; its integral is \(N\). A conventional spherical average of a density \(\rho\) is
\[
\widetilde\rho(r)=\frac1{4\pi}\int_{S^2}\rho(r\omega)\,\dd\omega,\qquad r\geq0.
\]
This is an average of density, not the radial shell distribution \(4\pi r^2\widetilde\rho(r)\). The source denotes its spherically averaged density by \(\rho_e\), but its original passage explicitly distinguishes a symmetrized physical density from a previously defined density whose definition is absent. The above formulas therefore explain standard conventions without silently identifying those two source densities.

The two retained questions are: determine the precise regularity of \(\rho_e\) near \(r=0\), and decide whether it is monotonically decreasing on its full radial domain. On the nonincreasing reading, the second target is
\[
\rho_e(r_2)\leq\rho_e(r_1)\qquad(0\leq r_1<r_2<\infty),
\]
for every ground state in the intended atomic class. The exact density convention and class must be fixed first. The question does not specify strict monotonicity or a particular regularity class. One-sided regularity of \(r\mapsto\rho_e(r)\) at zero and regularity of the radial function \(x\mapsto\rho_e(|x|)\) at the origin are distinct and should not be conflated.

The original passage also asks about positivity of \(\rho\) in \(\R^3\), then announces strict positivity of a spherically averaged ground-state density and an explicit lower bound at zero, using the symmetrized physical density. No bound or full hypotheses are reproduced. The corrected statement removes that announced result and retains the regularity and global-monotonicity questions. The supplied assessment checked 2026-08-21 reports positivity and local or asymptotic regularity results but not global monotonicity; it does not supply the exact remaining regularity target.
\subsection{Short English statement}
Determine the behavior of the atomic ground-state spherical density near the nucleus and whether it decreases at all radii. The physical versus previously defined density, symmetry sector, and precise regularity target must be recovered; the announced positivity result is background.
\subsection{Sources}
\begin{itemize}
\item Local UnsolvedMath registry, permanent ID 30000153, catalogue entry, \emph{Regularity and Monotonicity of Atomic Ground-State Densities}.
\item \emph{Spectral Analysis of Partial Differential Equations} (2005), Oberwolfach Reports, pp. 2839--2912, DOI \url{https://doi.org/10.4171/owr/2004/53}.
\item Supplied reference labeled ``current EMS report text'': \url{https://ems.press/content/serial-article-files/45974?nt=1}.
\item Original passage supplied in the registry, including the symmetry-subspace and physical-density qualifications; supplied literature triage checked 2026-08-21.
\end{itemize}
\end{document}
